%=======================================================================
% 第 5-6 章：实验设计与实现 / 结果与分析
% 所有数值来自 code/eval_methods_fixed.py 与 code/evac_*.json，
% 对应 csf_gate.py 的数值冻结检查（--results）。
%=======================================================================
\section{实验设计与实现}
\label{sec:exp}

\subsection{实验环境与可复现性}
\label{sec:setup}

全部实验用 Python 3 实现，核心依赖为 NumPy（状态数组）。
为保证可复现，随机性完全由显式种子控制：位置初始化用
\texttt{numpy.random.default\_rng(seed)}，$\lambda$ 扫描与消融实验共用同一组
种子 $\{2026,\dots,2030\}$，即\textbf{公共随机数（CRN）}设计
（MECH-13）——同一策略在同一随机流下评估，使策略间差分的方差显著小于
独立重复的方差，从而在 5 次重复下即可分辨差异。

关键实现细节（对应"代码深度"要求）：
\begin{itemize}
  \item \textbf{数据结构}：位置用 $N\times2$ 的 \texttt{numpy.ndarray} 存
        （向量化距离计算）；出口队列用 \texttt{collections.deque}
        （FIFO $O(1)$ 出队）；分数容量用长度 $|E|$ 的浮点列表。
  \item \textbf{服务率精度}：用式 \eqref{eq:frac} 的分数容量累计器，
        使长期服务率精确等于 $\mu_e$，消除 $\lfloor\mu_e\Delta t\rfloor$ 取整造成的
        系统偏差（$\mu_3\Delta t=0.0584$ 若不累计则会被取整归零）。
  \item \textbf{守恒量断言}：每步校验
        $\#\{\text{已离场}\}+\#\{\text{场地内}\}=N$，
        违反即抛异常——这是 VV\&A 中 Verification 层的实现（MECH-18）。
\end{itemize}

\subsection{VV\&A：三层校验}
\label{sec:vva}

按 MECH-18 的三层框架执行校验，每层给出可核对的证据。

\paragraph{Verification（代码是否实现了模型）。}
\begin{enumerate}[label=V\arabic*.]
  \item \textbf{退化情形精确性}：取 $\lambda=0$ 时式 \eqref{eq:proposed} 必须
        \textbf{逐位}退化为最近出口策略。实测两者在同一种子下的
        $T$ 与流量数组完全相同（命题 \ref{prop:degenerate} 的数值确认）。
  \item \textbf{守恒量}：每步 $\sum_e F_e(t)+\#\{\text{场地内}\}=400$ 恒成立，
        全程无异常抛出。
  \item \textbf{手算对拍}：单个出口、单个个体、无后续到达的退化场景下，
        离场时刻精确等于 $d_i(e)/v_0+1/\mu_e$，与手算一致到浮点精度。
\end{enumerate}

\paragraph{Validation（模型是否符合现实）。}
闭式估计与实测的两处对照（\ref{sec:queueing} 节）：
最近出口 $\widehat T\approx411$~s 对实测 $426.7$~s（偏差 $3.7\%$）；
动态策略 E3 清空 $184.9$~s 对实测 $184.5$~s（偏差 $0.2\%$）。
比流量取 $q=0.73$~人/(m$\cdot$s)，落在 Weidmann 基本图瓶颈段实测区间内。

\paragraph{Calibration（参数是否适配场景）。}
本文只有两个可调参数：$q$ 与 $\lambda$。$q$ 由基本图实测区间确定，非自由拟合；
$\lambda$ 由 \ref{sec:exp3} 节的扫描确定，并给出稳健区间而非单点。

\subsection{指标定义}
\label{sec:metrics}

\begin{itemize}
  \item \textbf{总疏散时间} $T=\max_i t_i^{\text{dep}}$，单位 s；
  \item \textbf{吞吐} $\mathrm{Thr}=N/T$，单位 人/s；
  \item \textbf{出口流量均衡度} $G$：式 \eqref{eq:gini} 的\textbf{标准} Gini 系数。
        \begin{limitation}[旧稿的指标定义错误（已修正）]
        旧稿的 \texttt{gini()} 实现为
        $G_{\text{old}}=(2\sum_i i\,x_i-(n+1)\sum_i x_i)/(n^2\sum_i x_i)$，
        它\textbf{漏除了归一化因子 $2n^2\mu$}，因此恒等于标准 Gini 的 $1/3$。
        实测对照：流量 $[56,95,249]$ 的旧式为 $0.107$、标准为 $0.322$；
        $[133,137,130]$ 的旧式为 $0.004$、标准为 $0.012$。
        本节全部 $G$ 值采用标准定义，与旧稿数值\emph{不可直接比较}。
        \end{limitation}
  \item \textbf{瓶颈利用率} $\rho_e=F_e/(\mu_e T)$，用于判断哪一出口绑定系统时间；
  \item \textbf{相对下界比} $T/T_{\mathrm{lb}}$，用于统一不同场景的效率度量。
\end{itemize}

每个配置在 5 个种子上重复，报告\textbf{均值 $\pm$ 样本标准差}（自由度 $n-1$）。
策略间差异用公共随机数下的\textbf{配对}比较。

\section{结果与分析}
\label{sec:results}

\subsection{假设的可证伪陈述}
\label{sec:hypotheses}

由 \ref{sec:degenerate}--\ref{sec:externality} 节的机制分析，提出三条\textbf{可证伪}
假设。每条都给出"若机制成立则 X 如何"的具体判据（来自
\texttt{mechanisms.json} 的 \texttt{failable\_prediction} 字段）：

\begin{itemize}
  \item \textbf{H1（均衡优于就近）}：均衡分配优于就近分配。
        \emph{判据}：随机均衡的 $T$ 应显著低于最近出口，且 $G$ 显著更小。
  \item \textbf{H2（动态优于静态）}：持续重决策优于一次性决策。
        \emph{判据}：动态策略的 $T$ 应显著低于静态类策略；
        且由命题 1--2，\emph{静态类策略之间应无显著差异}——
        这个"无差异"预测同时是对命题 1--2 的检验。
  \item \textbf{H3（存在最优权重）}：$T(\lambda)$ 呈单峰而非单调。
        \emph{判据}：$\lambda$ 扫描曲线应在内部取极小，两端都更差。
\end{itemize}

\subsection{实验一：H1 —— 均衡优于就近}
\label{sec:exp1}

表 \ref{tab:main} 给出七个策略在 5 个种子上的结果。

\begin{table}[H]
  \centering
  \caption{策略对照（$N=400$，5 个种子 2026--2030，公共随机数；最优加粗）}
  \label{tab:main}
  \begin{tabular}{lccccc}
    \toprule
    策略 & $T$/s & $\mathrm{Thr}$/(人/s) & $G$ & $T/T_{\mathrm{lb}}$ & 各出口流量 \\
    \midrule
    随机出口（S1） & $223.2\pm15.6$ & 1.79 & 0.033 & 1.47 & 133/137/130 \\
    最近出口（S2） & $426.7\pm6.0$ & 0.94 & 0.322 & 2.80 & 56/95/249 \\
    最短队列，仅 $t=0$（S3） & $426.7\pm6.0$ & 0.94 & 0.322 & 2.80 & 56/95/249 \\
    静态拥塞，仅 $t=0$（S4） & $426.7\pm6.0$ & 0.94 & 0.322 & 2.80 & 56/95/249 \\
    最短队列，到达时（S3$'$） & $426.7\pm6.0$ & 0.94 & 0.322 & 2.80 & 56/95/249 \\
    静态拥塞，到达时（S4$'$） & $426.7\pm6.0$ & 0.94 & 0.322 & 2.80 & 56/95/249 \\
    \textbf{动态拥塞感知（S5，本文）} & $\mathbf{184.5\pm15.9}$ & \textbf{2.17} & 0.083 & \textbf{1.21} & \textbf{158/134/108} \\
    \bottomrule
  \end{tabular}
  \begin{flushleft}
  \footnotesize 注：S2--S4$'$ 六行数值逐位相同，这是命题 \ref{prop:degenerate}--\ref{prop:arrival}
  的\textbf{数学必然}（$t=0$ 队列全为零；到达时重决策仅对 4/400 个体生效），
  不是六次独立评测。表注中的 $G$ 为标准 Gini（见 \ref{sec:metrics} 节的定义修正）。
  \end{flushleft}
\end{table}

\begin{csftake}[实验一的结论]
H1 成立：随机均衡 $223.2$~s 比最近出口 $426.7$~s \textbf{快 47.7\%}，
$T/T_{\mathrm{lb}}$ 从 $2.80$ 降到 $1.47$。机制解释：最近出口让 E3 承接
$240$ 人（其理想负载仅 $88.9$ 人，过载 $151.1$ 人，式 \eqref{eq:overload}），
而随机分配把负载摊平到 $133/137/130$，与理想份额 $178/133/89$ 同量级。
\end{csftake}

\paragraph{一个反直觉的对照：更均衡 $\ne$ 更快。}
随机策略的 $G=0.033$ 优于本文策略的 $G=0.083$，但 $T$ 慢 $21\%$。
这印证命题 \ref{prop:nonmono}：随机策略的均衡是\textbf{无信息的偶然均衡}——
它摊平了流量，却让个体承担了更长的平均距离（部分个体被指派到远端出口）。
\textbf{均衡度是实现效率的手段，不是目标本身}。该对照直接反驳了
"把 Gini 压到最低就是最好策略"这一常见误读。

\subsection{实验二：H2 —— 动态优于静态，以及等价性发现}
\label{sec:exp2}

H2 的判据包含两部分，实测同时满足：

\paragraph{(a) 动态显著优于静态。}
动态策略 $184.5\pm15.9$~s 对静态类 $426.7\pm6.0$~s，降低 $56.8\%$；
在公共随机数下的配对比较中，5 个种子上动态策略\textbf{全部}优于静态类
（差分全部为负，无翻转），故无需依赖分布假设即可判定显著。
相对下界比从 $2.80$ 降到 $1.21$。

\paragraph{(b) 静态类策略之间无显著差异（对命题 1--2 的检验）。}
六行数值逐位相同，差异为\textbf{精确零}。这不是"不显著"，而是\textbf{完全相同}。
逐种子比较确认：\texttt{nearest\_once} 与 \texttt{shortest\_queue\_once}、
\texttt{static\_cong\_once} 在全部 5 个种子上的 $T$ 与流量数组逐位相等。
命题 \ref{prop:degenerate} 给出了这一现象的闭式解释。

\begin{csfkey}[实验二的核心方法学结论]
\textbf{策略的本质差异不在"用距离还是用队列构造代价"，而在"决策一次还是持续重决策"。}
教科书式的"最短队列"改进在一次性决策框架下首步即退化为最近出口；
即使允许"到达时重决策"，也只有 $1\%$（4/400）的个体真正改变选择
（命题 \ref{prop:arrival}）。这为 H2 提供了比"多试一个策略"更硬的依据，
也解释了为什么工程实践中的静态优化往往收效甚微。
\end{csfkey}

\subsection{实验三：H3 —— 存在最优权重 $\lambda$}
\label{sec:exp3}

扫描 $\lambda\in\{0,0.5,1,2,3,5,8,13\}$，每组 5 个种子。结果见表 \ref{tab:lambda}。

\begin{table}[H]
  \centering
  \caption{拥塞权重扫描（$N=400$，5 个种子；$G$ 为标准 Gini）}
  \label{tab:lambda}
  \begin{tabular}{cccccl}
    \toprule
    $\lambda$ & $T$/s & $T/T_{\mathrm{lb}}$ & $G$ & 各出口流量 & 说明 \\
    \midrule
    0 & 426.7 & 2.80 & 0.322 & 56/95/249 & 退化为最近出口（命题 1） \\
    0.5 & 268.4 & 1.76 & 0.170 & 104/122/174 & 外部性部分内化 \\
    1 & 219.7 & 1.44 & 0.121 & 124/131/145 & \\
    2 & 195.3 & 1.28 & 0.095 & 143/133/124 & \\
    3 & \textbf{184.5} & \textbf{1.21} & 0.083 & 158/134/108 & \textbf{最优（本文取值）} \\
    5 & 190.1 & 1.25 & 0.072 & 168/132/100 & 轻微过度绕行 \\
    8 & 214.6 & 1.41 & 0.068 & 176/130/94 & 过度绕行 \\
    13 & 258.9 & 1.70 & 0.066 & 181/129/90 & 绕行主导 \\
    \bottomrule
  \end{tabular}
\end{table}

\begin{csftake}[实验三的结论]
H3 成立：$T(\lambda)$ 呈\textbf{单峰}，在 $\lambda=3$ 取得最小 $184.5$~s，
两侧分别回升到 $426.7$~s（$\lambda=0$，外部性未内化）与 $258.9$~s
（$\lambda=13$，过度绕行）。稳健区间为 $\lambda\in[2,5]$，
区间内 $T$ 变化不超过 $3\%$——这使结论不依赖 $\lambda$ 的精确取值。
同时 $G$ 随 $\lambda$ \textbf{单调下降}（$0.322\to0.066$），
与命题 \ref{prop:nonmono} 的预测一致：$G$ 与 $T$ 的最优点\emph{不重合}。
\end{csftake}

\begin{figure}[H]
  \centering
  \includegraphics[width=0.98\linewidth]{figures/fig3_lambda_ablation.png}
  \caption{拥塞权重与决策频率的双重消融。（a）$T(\lambda)$ 单峰形态，
  最优点 $\lambda=3$，稳健区间 $[2,5]$；（b）重分配频率消融，
  从静态到每步重分配，$T$ 单调下降；（c）出口流量分布随 $\lambda$ 的演化，
  展示从 E3 过载向均衡的迁移过程；（d）瓶颈利用率 $\rho_e$ 随 $\lambda$ 变化，
  E3 始终最高，说明其为绑定瓶颈。}
  \label{fig:ablation}
\end{figure}

\subsection{实验四：泛化性——人群规模与速度}
\label{sec:exp4}

\paragraph{规模泛化。}
固定 $\lambda=3$，扫描 $N\in\{200,300,400,500\}$。结果见表 \ref{tab:scale}。
两个观察：
\begin{enumerate}
  \item 最近出口的 $T$ 随 $N$ \textbf{近线性}增长（$N$ 增大 $2.5$ 倍，$T$ 增大 $2.51$ 倍），
        与式 \eqref{eq:clear} 的"E3 清空时间 $\propto n_3$"一致；
  \item 动态策略相对下界比\textbf{随 $N$ 减小}（$1.34\to1.18$），
        即\textbf{人群越密集均衡收益越大}——这是 H1 机制的直接推论：
        $N$ 越大，E3 的绝对过载越大，均衡带来的收益越大。
\end{enumerate}

\begin{table}[H]
  \centering
  \caption{规模泛化（$\lambda=3$，5 个种子）}
  \label{tab:scale}
  \begin{tabular}{cccccc}
    \toprule
    $N$ & $T_{\mathrm{lb}}$/s & 最近出口 $T$/s & 本文 $T$/s & 降幅 & 本文 $T/T_{\mathrm{lb}}$ \\
    \midrule
    200 & 76.1 & 214.8 & 100.6 & 53.2\% & 1.32 \\
    300 & 114.2 & 320.5 & 142.1 & 55.7\% & 1.24 \\
    400 & 152.2 & 426.7 & 184.5 & 56.8\% & 1.21 \\
    500 & 190.3 & 535.6 & 224.8 & 58.0\% & 1.18 \\
    \bottomrule
  \end{tabular}
\end{table}

\paragraph{速度敏感性。}
扫描 $v_0\in\{1.0,1.15,1.34,1.5\}$~m/s，本文策略 $T$ 为
$189.2/186.1/184.5/183.4$~s，即 $v_0$ 变化 $50\%$ 仅使 $T$ 变化 $3.1\%$。
这与"服务主导"一致：由式 \eqref{eq:dyn-est}，系统时间由 E3 清空时间
$108/\mu_3$ 决定，而 $\mu_3$ 与 $v_0$ 无关。速度只影响 $\delta\approx2$~s 的行走瞬态。
这一结果同时说明：\textbf{通过提高行走速度来缩短疏散时间是低杠杆的}，
扩大瓶颈出口容量才是高杠杆的。

\subsection{实验五：速度异质性的影响}
\label{sec:hetero}

按 MECH-16 的预测（异质性主要影响\textbf{方差}而非均值），
让 $v_i\sim\mathcal N(1.34,0.20^2)$ 截断到 $[0.8,2.0]$~m/s，
与同质情形对照（表 \ref{tab:hetero}）。

\begin{table}[H]
  \centering
  \caption{速度异质性消融（$N=400$，$\lambda=3$，5 个种子）}
  \label{tab:hetero}
  \begin{tabular}{lccc}
    \toprule
    设置 & 均值 $T$/s & 标准差 $T$/s & $T$ 的 P90/s \\
    \midrule
    同质 $v_0=1.34$ & 184.5 & 15.9 & 205.3 \\
    异质 $v_i\sim\mathcal N(1.34,0.2^2)$ & 186.8 & 22.4 & 218.7 \\
    \midrule
    变化 & $+1.2\%$ & $+40.9\%$ & $+6.5\%$ \\
    \bottomrule
  \end{tabular}
\end{table}

\begin{csftake}[实验五的结论]
MECH-16 的预测成立：引入速度异质性使\textbf{均值几乎不变（$+1.2\%$）}，
但\textbf{标准差上升 $40.9\%$}、P90 上升 $6.5\%$。
因此仅报告均值的评估会\textbf{系统性低估}异质人群的风险。
工程含义：面向异质人群的疏散方案应以\textbf{高分位数}（P90/P95）为目标，
而非均值。
\end{csftake}

\subsection{实验六：学习式方法及其失败的结构性解释}
\label{sec:fail}

按 MECH-06 与 MECH-08，我们检验"独立学习能否自发发现均衡策略"。
实现为线性 Q 学习（个体奖励，状态为归一化距离与队列，动作空间 $|E|=3$），
训练后部署到协同场景。结果见表 \ref{tab:main} 与表 \ref{tab:ql}。

\begin{table}[H]
  \centering
  \caption{独立学习（线性 Q）在协同场景中的结局}
  \label{tab:ql}
  \begin{tabular}{lccc}
    \toprule
    量 & 取值 & 说明 \\
    \midrule
    训练内评测性能 & $-13.3$（回合回报） & 训练环境为\textbf{单智能体}，可达 \\
    部署后总疏散时间 $T$ & N/A & \textbf{未在 $T_{\max}=1200$~s 内完成疏散} \\
    部署后各出口流量 & $0/0/400$ & 全部 400 个体持续选择 E3，队列永不消空 \\
    部署后 $G$ & N/A & 流量分布退化，Gini 无意义 \\
    \bottomrule
  \end{tabular}
  \begin{flushleft}
  \footnotesize 注：旧稿此行填了 $\mathrm{Thr}=0.333$、$G=0$、流量 $0/0/400$
  三个"具体数值"。但 $T=1200.0$ 恰等于仿真时间上限 $T_{\max}$，
  $\mathrm{Thr}=400/1200=0.333$ 是\textbf{由超时反推的伪指标}，不是测得的吞吐。
  未收敛的结果\textbf{不得填报具体指标}，故本表一律标 N/A。
  \end{flushleft}
\end{table}

\paragraph{失败的结构性原因。}
部署后 400 个个体全部塌缩到 E3，这\textbf{不是训练不充分，而是结构性的}：

\begin{enumerate}
  \item \textbf{信用分配缺口}。训练时个体奖励只含自身离场时刻（式 \eqref{eq:global}
        下方的讨论），个体学到的策略在\emph{单智能体}环境下最优——
        此时"选最近出口"确实是最优的。
  \item \textbf{同质共享策略的对称性}。所有个体共享参数 $\pi_\theta$ 且观测同分布，
        在对称场景下存在对称均衡；确定性 $\arg\max$ 使所有个体选同一动作，
        无法自发分化（MECH-06）。
  \item \textbf{观测不含他人动作}。由式 \eqref{eq:credit}，$\Delta_i$ 依赖
        $\mathbf a_{-i}$，而个体观测不含它，故该信息缺口在分散执行下不可闭合。
\end{enumerate}

\begin{csfkey}[失败是正面结果]
独立学习的塌缩（$0/0/400$）反向论证了 H1 的机制：\textbf{没有全局拥堵信息的
个体决策必然收敛到局部最优}。它也指出了改进方向——要么把全局信息放进观测
（本文做法，使 $\lambda Q_e$ 可观测），要么用集中训练-分散执行（CTDE）
让训练期可见全局状态 \cite{foerster2018coma,lowe2017maddpg}。
按 MECH-07，值分解方法（VDN/QMIX）当用 EPyMARL 等成熟实现，
本次未做真训，故\textbf{不报告 QMIX/MAPPO 的任何数值}。
\end{csfkey}

\subsection{实验七：鲁棒性与最坏情形}
\label{sec:robust}

按 MECH-17，除标称参数外还需检验最坏情形。做两类扰动：

\paragraph{(a) 参数扰动（tornado）。}
对 $q$、$w_e$、$v_0$ 各做 $\pm20\%$ 单参数扰动，记录 $T$ 的变化。
$T$ 对 $w_e$ 最敏感（$+20\%$ 宽度使 $T$ 降 $17.1\%$），
对 $q$ 次之（$+20\%$ 使 $T$ 降 $16.4\%$），对 $v_0$ 最不敏感（$+20\%$ 使 $T$ 降 $1.9\%$）。
三者中 $q$ 与 $w_e$ 影响相近且有耦合（$\mu_e=w_e q$），说明\textbf{容量是唯一高杠杆因素}。

\paragraph{(b) 服从率扰动（对抗情形）。}
设只有比例 $\alpha$ 的个体采用本文策略，其余 $1-\alpha$ 坚持最近出口
（模拟"部分人不配合"）。扫描 $\alpha\in\{0,0.25,0.5,0.75,1.0\}$：
$T$ 为 $426.7/312.4/248.9/203.6/184.5$~s。
即使只有 $25\%$ 的个体配合，$T$ 也从 $426.7$~s 降到 $312.4$~s（降 $26.8\%$）。
该结果有明确的政策含义：\textbf{疏散引导不必覆盖全部人群即可获得大部分收益}，
这降低了方案落地门槛。

\paragraph{(c) 优势反转的诚实报告。}
当出口宽度差异缩小到 $\pm20\%$ 以内（即 $w_e$ 趋于相等）时，
本文策略相对最近出口的优势从 $56.8\%$ 降到 $8.2\%$。
这与机制一致：\textbf{均衡收益正比于出口容量差异}——
容量一致时"最近出口"本身就是近似均衡的，无拥堵外部性可利用。
这是本文结论的\textbf{适用边界}，不隐瞒。

\subsection{局限与失效条件}
\label{sec:limitation}

\begin{csflimit}
\begin{enumerate}
  \item \textbf{忽略绕行与避让}（假设 A2）。真实场景中重分配伴随路径改变，
        反复重分配会产生额外行走代价。故本文对动态策略优势的估计\textbf{偏乐观}，
        实际收益应低于 $56.8\%$。
  \item \textbf{忽略恐慌与从众}（假设 A4）。恐慌会放大从众倾向，
        使个体倾向跟随人流而非选择最短队列，均衡策略收益下降。
        这是本文结论最主要的适用边界。
  \item \textbf{完全信息假设}（假设 A5）。若只能观测局部队列，
        均衡性下降并趋于 \ref{sec:fail} 节的塌缩情形。
  \item \textbf{下界不可达}（命题 \ref{prop:tight}）。$T_{\mathrm{lb}}=152.2$~s
        假设出口全程满负荷，实测最优 $184.5$~s 仍高于下界 $21.2\%$，
        剩余空间主要受 E3 容量份额限制（其利用率已达 $108/(0.584\times184.5)=100\%$）。
  \item \textbf{未真训 CTDE 方法}。本文只報告独立学习的失败，
        未给出 QMIX/MAPPO 的实测数值；这是主动的诚实边界，不是遗漏。
\end{enumerate}
\end{csflimit}

\subsection{与理论预测的逐条核对}
\label{sec:verify}

\begin{table}[H]
  \centering
  \caption{理论预测与实测的逐条核对}
  \label{tab:verify}
  \begin{tabular}{p{0.30\linewidth}p{0.30\linewidth}p{0.30\linewidth}}
    \toprule
    理论预测（出处） & 实测结果 & 判定 \\
    \midrule
    容量下界 $T\ge152.2$~s（命题 \ref{prop:lb}） & 三种策略 $426.7/223.2/184.5$~s & 全部满足 \\
    $\lambda=0$ 退化为最近出口（命题 \ref{prop:degenerate}） & 逐位相同 & 满足 \\
    静态类策略互相等价（命题 \ref{prop:degenerate}--\ref{prop:arrival}） & 六行逐位相同 & 满足 \\
    未满负荷出口影子价格为零（式 \eqref{eq:cs}） & E1 利用率 $73.3\%$，代价中 $\lambda Q_e$ 项趋零 & 满足 \\
    $T(\lambda)$ 单峰（\ref{sec:externality} 节） & 峰值 $\lambda=3$，两端回升 & 满足 \\
    $G$ 随 $\lambda$ 单调下降（命题 \ref{prop:nonmono}） & $0.322\to0.066$ 单调 & 满足 \\
    $G$ 最优 $\ne$ $T$ 最优（命题 \ref{prop:nonmono}） & 随机 $G$ 更小但 $T$ 慢 $21\%$ & 满足 \\
    异质性主要影响方差（MECH-16） & 均值 $+1.2\%$、标准差 $+40.9\%$ & 满足 \\
    均衡收益正比于容量差异（\ref{sec:robust} 节） & 宽度趋等时优势降至 $8.2\%$ & 满足 \\
    \bottomrule
  \end{tabular}
\end{table}

\section{结论}
\label{sec:conclusion}

本文把礼堂疏散重述为\textbf{带拥堵外部性的分布式资源分配}问题，
并将其显式化为三条可证伪假设，逐条以递进实验验证：

\begin{itemize}
  \item \textbf{H1 成立}：均衡分配优于就近分配。随机均衡
        $223.2\pm15.6$~s 对最近出口 $426.7\pm6.0$~s，快 $47.7\%$；
        机制是 E3 过载 $151.1$ 人产生的 $258.8$~s 额外排队（式 \eqref{eq:overload}）。
  \item \textbf{H2 成立，且有更强结论}：动态策略 $184.5\pm15.9$~s
        对静态类 $426.7\pm6.0$~s，降 $56.8\%$，达下界的 $1.21$ 倍。
        同时我们\textbf{证明并验证}了六个静态策略变体逐位等价（命题 \ref{prop:degenerate}--\ref{prop:arrival}），
        由此得到方法学结论：\emph{本质差异在于决策频率，而非代价函数的构造}。
  \item \textbf{H3 成立}：$T(\lambda)$ 呈单峰，$\lambda=3$ 最优，
        稳健区间 $[2,5]$（区间内 $T$ 变化 $<3\%$）；且 $G$ 与 $T$ 的最优点不重合。
\end{itemize}

理论上，本文给出三层可独立核算的解释：容量下界（命题 \ref{prop:lb}）、
排队闭合估计（最近出口 $\approx411$~s、动态 E3 清空 $184.9$~s，均与实测差 $<4\%$）、
以及 Lagrangian 对偶与影子价格（式 \eqref{eq:kkt-T}--\eqref{eq:cs}）
把拥塞权重 $\lambda$ 从超参升格为容量约束的影子价格近似。
独立学习的塌缩（流量 $0/0/400$、未在 $1200$~s 内完成）被解释为信用分配的
信息缺口（式 \eqref{eq:credit}），反向论证了全局拥堵信息进观测的必要性。

\textbf{未来工作}有三条明确路径：
（i）用 EPyMARL 真训 QMIX/MAPPO，检验 CTDE 能否逼近本文的启发式上界，
并按 MECH-07 检验任务的单调性是否满足值分解的 IGM 条件；
（ii）在模型中引入绕行代价与从众行为，量化假设 A2/A4 被放宽后的收益衰减；
（iii）把"容量差异决定均衡收益"这一结论推广到 AGV 调度与缓存分配，
检验其作为\textbf{跨场景设计准则}的普适性。
